% Capitolo 4 — Sfide strutturali: pressione demografica e carriere discontinue
% Ultima modifica: 2026-05-28
% Fonti principali usate: MEF_RGS2025, MEF_RGS2025_NdA, FrancoTommasino2020,
%   Mazzaferro2023, Raitano2020, INPS2025, Censis_Confcooperative2026,
%   OCPI2026, OCSE2023, ISTAT_Natalita2024, EUROSTAT_NEET2024,
%   BosiGuerra2022, Bosi2018

%% ───────────────────────────────────────────────────────────────────
\chapter{Sfide strutturali: pressione demografica e carriere discontinue}
\label{ch:sfide}
%% ───────────────────────────────────────────────────────────────────

% --- §4.1 Dal sistema all'individuo ------------------------------------
\noindent
Il capitolo precedente ha costruito una mappa delle riforme possibili e ne ha
misurato il costo. La mappa è utile, ma non risponde alla domanda che un
lavoratore di venticinque anni si pone nel 2025: non \emph{cosa dovrebbe fare
il sistema}, ma \emph{cosa otterrà il sistema}. La differenza non è retorica.
Le riforme percorribili, come mostrato, liberano al massimo dieci o quindici
punti di tasso di sostituzione per le platee più svantaggiate, a condizione
che vengano implementate, che il vincolo di bilancio regga e che la produttività
italiana recuperi una traiettoria che nel decennio 2015-2025 è stata dello
0,12\% annuo \cite{OCPI2026}. Tre condizioni che non si possono dare per certe.

\medskip
\noindent
Ciò che invece è certo, dentro i margini di qualunque scenario
macroeconomico credibile, è che i meccanismi automatici del sistema NDC
traducono la pressione demografica in riduzioni del tasso di sostituzione in
modo sistematico e misurabile. Questo capitolo quantifica quella
sistematicità. L'analisi si sviluppa su due piani. Il primo è aggregato:
la struttura demografica italiana e la sua traiettoria di medio-lungo
periodo. Il secondo è individuale: quattro profili di carriera tipici della
Generazione Z e le loro pensioni attese, calcolate con gli stessi parametri
del metodo NDC già introdotti nei capitoli precedenti. La distanza tra i
due piani è la distanza tra la sostenibilità del sistema e l'adeguatezza per
la persona specifica, che è il nucleo della domanda di ricerca di questa tesi.

% --- §4.2 Il quadro demografico ----------------------------------------
\section{La pressione demografica come vincolo strutturale}
\label{sec:demografia}

\noindent
Il tasso di fecondità totale italiano è sceso a 1,20 figli per donna nel
2023, il valore più basso mai registrato nel dopoguerra, contro un valore
di sostituzione di 2,1 \cite{ISTAT_Natalita2024}. Nella proiezione base della
Ragioneria Generale dello Stato, il recupero parziale a 1,36 entro il 2070
lascia il rapporto demografico di dipendenza degli anziani, calcolato come
rapporto tra popolazione con 65 anni o più e popolazione in età lavorativa
(15-64), salire dal 37,2\% del 2024 al 59,0\% del 2070 \cite{MEF_RGS2025}.
In altre parole, una platea di pensionati che tra quarantacinque anni varrà
sei decimi della platea dei lavoratori finanzierà ogni anno una quota di
spesa pensionistica che nessuna riduzione dei coefficienti di trasformazione
può compensare interamente sul piano del benessere individuale.

\medskip
\noindent
La letteratura di scienza delle finanze formalizza questa dipendenza
attraverso il teorema di Samuelson-Aaron, che identifica il rendimento
implicito del sistema a ripartizione con la somma del tasso di crescita
dell'occupazione e del tasso di crescita dei salari reali:

\begin{equation}
r_{\text{PAYG}} = n + g_w
\label{eq:aaron}
\end{equation}

\noindent dove $n$ è il tasso di variazione dell'occupazione totale e
$g_w$ il tasso di crescita del salario medio reale \cite{Bosi2018}.
I due componenti sono entrambi sotto pressione per l'Italia.
L'occupazione nativa è in contrazione strutturale nella proiezione
demografica di base, mentre il contributo migratorio parzialmente
compensa senza rovesciare la tendenza \cite{MEF_RGS2025_NdA}.
La produttività del lavoro, che nei sistemi NDC determina il tasso di
capitalizzazione nozionale, ha registrato una crescita media dello
0,12\% annuo nel decennio 2015-2025 \cite{OCPI2026}; la proiezione
RGS di base incorporate un recupero all'1,04\% annuo nel lungo periodo,
ipotesi per cui lo stesso MEF segnala i gradi di incertezza quando
scrive che le proiezioni si fondano su un tasso di crescita
\guillemotleft{}lontano dalla dinamica recente\guillemotright{}
\cite{MEF_RGS2025_NdA}.

\medskip
\noindent
L'implicazione per un lavoratore che entra nel mercato del lavoro nel 2025
è diretta. Il rendimento implicito che il sistema PAYG gli riconosce sul
lungo periodo si colloca, nelle ipotesi prudenziali, nell'ordine dello
0,5-1,5\% reale annuo. Qualunque asset finanziario diversificato su un
orizzonte di quaranta anni restituisce storicamente un rendimento reale
di tre-cinque punti percentuali superiore \cite{OCSE2023}. La perdita
relativa rispetto a un ipotetico sistema a capitalizzazione non è più
una questione teorica: è la misura quantitativa del rischio demografico
che la generazione attuale di lavoratori porta in dote.

% --- §4.3 La Generazione Z nel mercato del lavoro ----------------------
\section{La Generazione Z e il mercato del lavoro italiano}
\label{sec:genz}

\noindent
I profili di carriera della Generazione Z, intesi come i nati tra il 1997
e il 2012, presentano caratteristiche strutturalmente diverse da quelle
delle coorti precedenti, non per scelta ma per composizione del mercato
del lavoro che si trovano di fronte. Il tasso di NEET, giovani tra i 15
e i 29 anni né occupati né in istruzione né in formazione, ha raggiunto
il 16,1\% in Italia nel 2023, il valore più alto dell'Unione europea
dopo la Romania \cite{EUROSTAT_NEET2024}. Tra gli occupati nella fascia
25-34 anni, il 27\% dei dipendenti è inquadrato con contratto a tempo
determinato, contro una media UE del 18\% \cite{INPS2025}. Questi numeri
non descrivono una transizione generazionale: descrivono il montante
contributivo che verrà accumulato, o non accumulato, nei prossimi
quarant'anni.

\medskip
\noindent
Il funzionamento del metodo contributivo nozionale trasforma ogni
discontinuità di carriera in una riduzione permanente e non recuperabile
del montante. Nel sistema retributivo, la pensione era funzione degli
ultimi anni di carriera: un lavoratore che aveva avuto un percorso
accidentato poteva compensare parzialmente con retribuzioni crescenti
verso la fine della vita attiva. Nel metodo NDC quella funzione di
smoothing non esiste. Il contributo non versato all'anno $t$ non viene
mai capitalizzato, perché l'anno $t$ è definitivamente passato. Una
pausa di cinque anni a inizio carriera, tipica dei percorsi di
transizione tra istruzione e lavoro stabile, riduce il montante finale
in misura equivalente a perdere i cinque anni di capitalizzazione più
lunga dell'intera vita lavorativa, ovvero quelli che avrebbero beneficiato
del fattore di crescita composto più alto. Il costo previdenziale di
un'interruzione precoce è sistematicamente superiore al costo di
un'interruzione tardiva.

\medskip
\noindent
A questo si aggiunge un effetto di composizione documentato dall'INPS:
i lavoratori con carriere discontinue tendono a concentrarsi nei settori
a bassa crescita salariale, il che riduce sia il numeratore (contributi
versati) sia il tasso di capitalizzazione del montante, che dipende dalla
crescita media del PIL nominale \cite{INPS2025}. Due canali di
compressione del montante che operano simultaneamente e si autorinfor\-zano
nel tempo.

% --- §4.4 Simulazioni su quattro profili di carriera -------------------
\section{Simulazioni: il costo previdenziale delle carriere discontinue}
\label{sec:simulazioni}

\noindent
Per rendere concrete le implicazioni appena esposte, si simulano quattro
profili di carriera rappresentativi della Generazione Z. Il punto di
partenza è la formula del montante nozionale già presentata nell'introduzione:

\begin{equation}
MC_j = c \cdot S \cdot \frac{(1+g)^{T_j} - 1}{g}
\label{eq:montante_profili}
\end{equation}

\noindent dove $c = 0{,}33$ è l'aliquota contributiva per il lavoro
dipendente, $S = 25.000$ euro il salario lordo annuo costante,
$g = 1{,}5\%$ il tasso di capitalizzazione nozionale (ipotesi prudenziale
coerente con lo scenario base RGS \cite{MEF_RGS2025_NdA}),
$T_j$ il numero di anni di contribuzione del profilo $j$. Tutti i profili
assumono il pensionamento a 67 anni nel 2067 e un coefficiente di
trasformazione del 5,0\% \cite{MEF_RGS2025}.
\todo{VERIFICA DATO: il coefficiente per il 2067 non è direttamente
stimato dal MEF-RGS2025 a quella data specifica; si usa la proiezione
tendenziale per difetto, coerente con l'andamento già previsto per il 2040
(5,2\%) e la speranza di vita al 2070 (+4,8 anni per gli uomini).}

\subsection*{Profilo A -- Dipendente stabile (quarant'anni di contribuzione)}

\noindent
Ingresso nel mercato del lavoro a 27 anni, carriera continua fino al
pensionamento. Quarant'anni di contribuzioni piena.

\[
MC_A = 8.250 \times \frac{(1{,}015)^{40} - 1}{0{,}015}
= 8.250 \times 54{,}27 \approx 447.700 \text{ euro}
\]

\[
P_A = 0{,}050 \times 447.700 \approx 22.385 \text{ euro/anno}
\quad (1.722 \text{ euro/mese})
\]

Il tasso di sostituzione lordo è $22.385 / 25.000 = 89{,}5\%$, valore
prossimo alla soglia di adeguatezza convenzionalmente fissata all'80\%
per i sistemi a prestazione definita.

\subsection*{Profilo B -- Gap contributivo iniziale (trentacinque anni)}

\noindent
Cinque anni di disoccupazione o lavoro irregolare tra i 27 e i 32 anni
(NEET, contratti a progetto non coperti da contribuzione piena, o
formazione post-laurea non retribuita). Trentacinque anni di
contribuzioni effettive.

\[
MC_B = 8.250 \times \frac{(1{,}015)^{35} - 1}{0{,}015}
= 8.250 \times 45{,}59 \approx 376.100 \text{ euro}
\]

\[
P_B = 0{,}050 \times 376.100 \approx 18.805 \text{ euro/anno}
\quad (1.446 \text{ euro/mese})
\]

Tasso di sostituzione lordo: $75{,}2\%$. La differenza rispetto al
profilo A è di 14,3 punti percentuali e di circa 3.580 euro annui
lordi, per effetto di cinque anni mancanti posizionati all'inizio della
carriera, dove il fattore di capitalizzazione composta opera per
quarant'anni contro i trentacinque del profilo stabile.

\subsection*{Profilo C -- Dieci anni a regime ridotto (caregiving)}

\noindent
Carriera di quarant'anni in cui i primi dieci si svolgono a metà orario
per carichi di cura (salario effettivo 12.500 euro, contribuzione
annua 4.125 euro). Tipico di lavoratrici con figli o con genitori
anziani a carico, che interrompono il lavoro a tempo pieno ma non
escono dal mercato.

Il montante si decompone in due periodi:

\[
\begin{aligned}
MC_{C,1}
  &= 4.125 \times \frac{(1{,}015)^{10} - 1}{0{,}015} \times (1{,}015)^{30} \\
  &= 4.125 \times 10{,}70 \times 1{,}563 \approx 69.010 \text{ euro}\\[6pt]
MC_{C,2}
  &= 8.250 \times \frac{(1{,}015)^{30} - 1}{0{,}015}
  = 8.250 \times 37{,}54 \approx 309.700 \text{ euro}
\end{aligned}
\]

\[
MC_C = 69.010 + 309.700 \approx 378.700 \text{ euro}
\]

\[
P_C = 0{,}050 \times 378.700 \approx 18.935 \text{ euro/anno}
\quad (1.456 \text{ euro/mese})
\]

Tasso di sostituzione lordo: $75{,}7\%$ (calcolato sul salario
effettivo medio di 18.750 euro). Rispetto al profilo A, la perdita
è di circa 3.450 euro annui lordi, distribuita su dieci anni di
contribuzione parziale invece che su cinque anni di assenza completa.
L'effetto finale è quasi identico al profilo B, confermando che
la discontinuità previdenziale si genera tanto dall'assenza quanto
dalla presenza ridotta.

\subsection*{Profilo D -- Autonomo con reddito sotto-dichiarato}

\noindent
Lavoratore autonomo per quarant'anni, con aliquota contributiva del
25\% (gestione separata INPS per autonomi) e reddito dichiarato di
18.750 euro a fronte di un reddito effettivo di 25.000 euro. La
contribuzione annua è 4.688 euro, il 56,8\% di quella del dipendente.

\[
MC_D = 4.688 \times 54{,}27 \approx 254.400 \text{ euro}
\]

\[
P_D = 0{,}050 \times 254.400 \approx 12.720 \text{ euro/anno}
\quad (978 \text{ euro/mese})
\]

Il tasso di sostituzione lordo rispetto al reddito effettivo (25.000
euro) è del $50{,}9\%$, tredici punti al di sotto della soglia di
adeguatezza dell'80\% e quasi quaranta punti sotto il profilo stabile.
La sotto-dichiarazione del reddito, fenomeno strutturale nel lavoro
autonomo italiano documentato dall'INPS \cite{INPS2025}, produce una
trappola previdenziale: l'autonomo che evade parzialmente l'imposta
oggi costruisce una povertà pensionistica per sé stesso domani.

\begin{table}[h!]
\centering
\small
\caption{Simulazioni su quattro profili Gen-Z: montante, pensione e
         tasso di sostituzione lordo}
\label{tab:profili}
\begin{tabularx}{\textwidth}{X r r r r}
\toprule
\textbf{Profilo} & \textbf{Anni} & \textbf{MC (euro)}
                 & \textbf{Pensione (euro/a)}
                 & \textbf{TS lordo} \\
\midrule
A -- Dipendente stabile       & 40 & 447.700 & 22.385 & 89,5\% \\
B -- Gap 5 anni (NEET/precario)& 35 & 376.100 & 18.805 & 75,2\% \\
C -- Regime ridotto 10 anni   & 40\,$^{a}$ & 378.700 & 18.935 & 75,7\% \\
D -- Autonomo sotto-dichiarato & 40 & 254.400 & 12.720 & 50,9\% \\
\bottomrule
\end{tabularx}
\par\smallskip
\footnotesize
Ipotesi comuni: salario lordo 25.000 euro costante, capitalizzazione
nozionale 1,5\%, coefficiente 2067 = 5,0\%, pensionamento a 67 anni.\\
$^{a}$ Di cui 10 anni a regime ridotto (50\%). Tasso di sostituzione
calcolato sul salario medio effettivo di 18.750 euro.\\
Fonti: elaborazioni dell'autore su parametri MEF-RGS \cite{MEF_RGS2025}
e INPS \cite{INPS2025}.
\end{table}

% --- §4.5 La regressività implicita del metodo NDC --------------------
\section{La regressività implicita: chi paga il prezzo della neutralità attuariale}
\label{sec:regressivita}

\noindent
La tabella~\ref{tab:profili} fotografa una distribuzione di esiti
che il metodo NDC, per costruzione, non considera ineguale. Ogni
lavoratore riceve in pensione il valore attualizzato di quanto ha
versato, modulato dalla longevità della propria coorte: non c'è
privilegi~o né discriminazione nel calcolo. La neutralità attuariale è
rispettata. Eppure i quattro profili mostrano tassi di sostituzione
compresi tra il 90\% e il 51\%, tutti partendo dallo stesso salario
lordo. La neutralità formale coesiste con una regressività sostanziale,
perché le discontinuità di carriera che producono questi esiti divergenti
non sono distribuite casualmente nella forza lavoro.

\medskip
\noindent
Franco e Tommasino identificano questa tensione come la contraddizione
fondamentale del sistema NDC italiano: il metodo garantisce equità
attuariale tra chi ha carriere omogenee, ma non è un meccanismo di
redistribuzione e non può esserlo per definizione
\cite{FrancoTommasino2020}. I lavoratori che accumulano carriere
più discontinue, giovani a bassa qualificazione, donne con carichi
di cura, autonomi in settori a basso reddito, abitanti del
Mezzogiorno, sono esattamente quelli per i quali il risparmio
previdenziale integrativo è meno accessibile. Il doppio svantaggio,
pensione pubblica bassa e capacità contributiva al pilastro
complementare limitata, non è una coincidenza: nasce dalla correlazione
tra posizione nel mercato del lavoro e struttura della carriera.

\medskip
\noindent
La dimensione di genere è il caso più documentato di questa
regressività. Le donne accumulano in media sette anni di interruzioni
di carriera legate al caregiving nel corso della vita lavorativa
\cite{INPS2025}.
\todo{VERIFICA DATO: il dato di 7 anni è approssimativo; l'INPS 2025
non fornisce questa stima in forma aggregata. Si verifica sul Rapporto
INPS 2025, sezione carriere femminili.}
Per le lavoratrici entrate nel regime contributivo puro dopo il 1995,
il divario di genere nel montante atteso è stimato tra il 23 e il 28\%
\cite{MEF_RGS2025}, contro un divario del 16-18\% per le coorti
che beneficiavano ancora di un regime misto retributivo-contributivo.
Il metodo NDC, che pure non distingue per sesso, amplia il gender
pension gap invece di ridurlo, perché non compensa le asimmetrie
nella distribuzione del lavoro di cura.

\medskip
\noindent
Michele Raitano quantifica la platea strutturalmente a rischio di
inadeguatezza: circa 1,2 milioni di nuove domande di pensione all'anno
entro il 2040 presenteranno un montante NDC insufficiente a raggiungere
la soglia minima di accesso al pensionamento anticipato, fissata in tre
volte l'assegno sociale dal 2030 \cite{Raitano2020}. Questi lavoratori
non riceveranno una pensione bassa: rischiano di non riceverne una
ordinaria per decenni. Il Profilo B e il Profilo C della tabella
precedente si collocano entrambi al di sopra di questa soglia, ma solo
a condizione che il salario e il tasso di capitalizzazione reggano
lungo tutto il percorso. Un'ulteriore riduzione del 15-20\% del
montante, prodotta da periodi di disoccupazione più lunghi o da
retribuzioni inferiori alla media, li spingerebbe sotto.

\medskip
\noindent
Carlo Mazzaferro ha documentato che il 92,6\% degli individui
pensionati tra il 1996 e il 2019 ha ricevuto in valore attuale più di
quanto versato \cite{Mazzaferro2023}. Il dato, già citato nel capitolo
precedente, acquista qui un'ulteriore valenza: quella generosità si è
tradotta in un sussidio implicito di 80 miliardi di euro a valori
presenti, finanziato per costruzione dalle coorti successive. La
generazione Z non eredita soltanto un sistema adeguato al suo
sostenimento: eredita il conto sospeso delle generazioni che l'hanno
preceduta, che il metodo NDC pone a carico dei loro montanti attraverso
un tasso di capitalizzazione nozionale agganciato alla crescita del PIL
nominale di un'economia che cresce poco.

\medskip
\noindent
Il dibattito tra gli esperti su questo punto non è privo di sfumature.
Tito Boeri e Pietro Garibaldi, analizzando la struttura intergenerazionale
del sistema pensionistico italiano, sostengono che la correzione della
regressività non possa venire dall'interno del primo pilastro senza
rompere l'equilibrio finanziario: essa richiede trasferimenti espliciti
dalla fiscalità generale verso i lavoratori più svantaggiati, nella forma
della pensione contributiva di garanzia o di contribuzione figurativa per
periodi di cura \cite{LavoceGiovani}. La proposta ha una logica interna
coerente, ma il suo costo è quello che il capitolo precedente ha stimato
tra tre e cinque miliardi annui, finanziamento che compete con ogni altra
voce di spesa pubblica in un paese con il 137\% di debito pubblico in
rapporto al PIL.

% --- §4.6 Conclusione critica ------------------------------------------
\section{Il rischio individuale senza strumenti individuali}
\label{sec:rischio_individuale}

\noindent
Le simulazioni della sezione~\ref{sec:simulazioni} dimostrano che
il sistema NDC, tecnicamente neutro, produce esiti radicalmente diversi
per profili di carriera la cui distribuzione nella forza lavoro
non è neutra. La variabile che separa il 89,5\% del Profilo A dal
50,9\% del Profilo D non è né il salario (uguale per entrambi) né il
numero di anni di lavoro (uguale o superiore nel Profilo D): è la
qualità dell'integrazione nel mercato formale, che il metodo NDC non
corregge ma amplifica.

\medskip
\noindent
Il punto che vale la pena tenere aperto, invece di chiudere con
una raccomandazione ordinata, è quello della distribuzione della
capacità di risposta. Il capitolo successivo presenterà gli strumenti
di pianificazione individuale disponibili: TFR al fondo pensione,
deducibilità fiscale dei contributi volontari, secondo pilastro
occupazionale, ETF e piani di accumulo come complemento. Questi
strumenti funzionano, e i dati della Ragioneria Generale mostrano
che un lavoratore del Profilo B che aderisce al fondo pensione sin
dall'inizio recupera fino a 7-10 punti di tasso di sostituzione
\cite{MEF_RGS2025}. Ma l'accesso a questi strumenti è condizionato
da un prerequisito che non riguarda la conoscenza finanziaria: è il
reddito disponibile. Un lavoratore con €25.000 annui e due figli a
carico non ha lo stesso margine di manovra di uno con €40.000 e
carriera stabile. Il pilastro complementare non è neutro rispetto
alla distribuzione del reddito: è efficiente per chi ha già qualcosa,
meno per chi non ce l'ha.

\medskip
\noindent
Questo non è un argomento contro il pilastro complementare. È un
argomento per costruirlo sapendo chi esclude strutturalmente, e per
non affidarsi a esso come unica risposta al problema
dell'adeguatezza. La generazione Z affronta un sistema che ha
trasferito su di lei il rischio demografico attraverso meccanismi
automatici, e al tempo stesso si trova in un mercato del lavoro che
rende quella risposta individuale sistematicamente più difficile di
quanto non fosse per i propri genitori. La distanza tra l'urgenza
della risposta e la capacità di darla è la vera sfida strutturale
di cui il titolo di questo capitolo porta il nome.
